\scanpage{15}
\uline{Proposition} 2 \quad \unsure{Soient|Soit} $A$ un anneau local artinien, $A'$ une $A$-algèbre finie, $A''=(A'\otimes_A A')/J$, où $J$ est un idéal d'équivalence, \deleted{On suppose que} \inserted{$\Delta=\operatorname{Ker}(A''\mathrel{\unsure{-|\longrightarrow}} A')$}. On suppose que \deleted{pour tt $x\in A'$}
\diagram{\begin{tikzcd}[ampersand replacement=\&]
A\arrow[r]\& A'\arrow[r,shift left=.5ex,"p_1"]\arrow[r,shift right=.5ex,"p_2"']\& A''
\end{tikzcd}}
est une suite \uline{exacte}, et que pour tt $x'\in A'\mathbin{-}A$, l'élément
\[
\delta(x')=p_2(x')-p_1(x')
\]
est un \uline{générateur} de l'idéal $\Delta$.
% slanted
\inserted{Supposons de plus les extensions résiduelles de $A'/A$ \uline{triviales}.}
% slanted
\margin{Sans l'hypothèse $\uparrow$, il est vrai que $\unsure{m|\mathfrak m}A'=m$, donc on est ramené à \unsure{étudier|discuter} le cas $A=k$}
Sous ces conditions, $\operatorname{long} A'/A\leq1$, et par suite (cf prop 1, (i)) on a $A''\simeq A'\otimes_A A'$.

Soit $m'=\unsure{\mathfrak r|r}(A')$, distinguons deux cas

1\textsuperscript{o}) $A'$ \uline{n'est pas \deleted{local},} \inserted{radiciel sur $A$} \deleted{Alors $\Delta$ n'est pas} \inserted{i.e. $\operatorname{Ker}(A'\otimes_A A'\longrightarrow A')$} n'est pas nilpotent. \uline{Comme $J$ est nilpotent}, il s'ensuit que $\operatorname{Ker}(A''\longrightarrow A')=\Delta$ n'est pas nilpotent. Mais si $x'\in\sqrt{\unsure{A|A'}}=m'$ \deleted{\unsure{c'est}} alors $\delta(x')$ est nilpotent, donc $\delta(x')A''$ est nilpotent, d'où $\delta(x')A''\ne\Delta$, d'où $x'\in A$ par hyp, d'où
\[
m'\subset A
\]
\inserted{et comme $\deleted{\unsure{mA'|mA}}m'\cap A\subset m\subset mA'\subset m'$, on a} d'où
\[
m'=m=mA'
\]
% archive sheet number: 14
